\section{Introduction}

Drug-target affinity (DTA) prediction links chemical structure to quantitative binding strength, enabling virtual screening, lead prioritization, and mechanism-aware repurposing \cite{vamathevan2019ai_drug,chen2018rise}. Deep learning models---DeepDTA \cite{ozturk2018deepdta}, GraphDTA \cite{nguyen2021graphdta}, MolTrans \cite{huang2021moltrans}, ConPlex \cite{singh2023conplex}, among others---have substantially improved benchmark accuracy by learning joint representations of proteins and small molecules from sequence, graph, or language-model features. Yet these gains have been driven almost entirely by in-distribution evaluation, where training and test data share the same proteins, compounds, and assay conventions. In practical drug discovery, a model is rarely applied to data drawn from its training distribution; it must score compounds against proteins from different assays, different families, or entirely different curated resources.

This gap is not hypothetical. DeepDTA achieves an AUROC of 0.898 on the held-out test split of Drug Target Commons (DTC) \cite{tang2014dtc}, yet drops to 0.534 when evaluated on the Davis kinase benchmark \cite{davis2011kinase} under verified 0\% drug overlap with training. The implication is that current models do not reliably learn portable determinants of binding; instead, they exploit dataset-specific regularities---recurring protein families, assay conventions, chemical neighborhoods---that vanish once evaluation moves outside the source benchmark \cite{pahikkala2015kronrls,mayr2018deeptox}.

We argue that this fragility reflects a formulation problem, not merely a representation problem. Binding evidence is inherently relational: if a drug is known to bind one protein strongly, that fact constrains which other proteins are likely to bind the same drug. Conventional DTA architectures compress each protein--drug pair into an isolated prediction, discarding this relational structure. What is missing is an explicit mechanism to reuse known interactions as transferable evidence when the query lies outside the training distribution \cite{pan2010transfer,zhuang2021transfer}.

This observation motivates a simple reformulation. Suppose protein~$A$ is a known strong binder of drug~$X$, and protein~$B$ is a query for which affinity to~$X$ is unknown. The known interaction $(A, X)$ establishes a reference: $X$ is compatible with at least one protein context, and the relevant question becomes whether $B$ shares the binding-relevant features of that context. If~$B$ resembles~$A$ in the aspects that govern interaction with~$X$, then the observed affinity between~$A$ and~$X$ should inform the prediction for~$B$. We call this \emph{anchor transfer}: using known strong binders as bridges that carry binding knowledge from observed interactions to unobserved ones.

Anchor transfer leads to a triplet formulation of DTA prediction. Instead of scoring a protein--drug pair $(q, d)$ alone, the model receives a triplet $(a, q, d)$, where the anchor~$a$ is a known strong binder of~$d$. The model is not asked to infer binding from scratch; it is asked whether the binding context established by the anchor supports affinity for the query. Figure~\ref{fig:architecture} provides an overview.

\begin{figure}[t]
    \centering
    \includegraphics[width=\textwidth]{fig1_architecture.png}
    \caption{Overview of the Anchor Transfer DTA architecture. Given a query protein and a drug, the model retrieves a known strong binder (anchor protein) of that drug. Both proteins are encoded with frozen ESM-2 embeddings projected to a shared 256-dimensional space. The drug is encoded with a parallel SMILES CNN. Three pairwise interaction MLPs model anchor--drug, query--drug, and anchor--query relationships. The fused 1536-dimensional representation feeds dual prediction heads for binding classification and pKi regression.}
    \label{fig:architecture}
\end{figure}

The architecture is deliberately simple. Proteins are encoded with frozen ESM-2 embeddings \cite{lin2023esm2}, drugs with a convolutional encoder over SMILES strings \cite{weininger1988smiles}, and the triplet is processed through three pairwise multilayer perceptrons that model anchor--drug, query--drug, and anchor--query interactions. This keeps the model lightweight while making the transfer pathway explicit.

Because anchor transfer targets generalization rather than benchmark optimization, we evaluate under conditions designed to expose cross-dataset failure. Using DTC as the source, we test on the Davis kinase benchmark \cite{davis2011kinase} (442 kinases, 68 drugs, zero drug overlap with training).

The results confirm that explicit relational transfer is more robust than isolated pair scoring. Crucially, anchors are retrieved at test time from the training set by Tanimoto chemical similarity after excluding canonical chemical duplicates---no oracle access or molecular leakage is involved. On Davis, V2-650M achieves per-protein CI of 0.642 and AUROC of 0.669, while DeepDTA falls to random. Stratifying by anchor binding strength reveals a complementary relationship between encoder capacity and anchor quality: the 650M encoder compensates for weak anchors, while strong anchors equalize the advantage of larger encoders. A cosine similarity ablation shows that 74\% of the model's predictive variance is orthogonal to raw protein similarity, confirming that it learns drug-specific binding patterns beyond nearest-neighbor retrieval.

This paper makes three contributions:
\begin{enumerate}
    \item An \emph{anchor transfer architecture} for DTA that conditions each query on a known strong binder of the same drug, using ESM-2 protein embeddings, a SMILES convolutional encoder, and triple pairwise interaction heads.
    \item A \emph{cross-dataset evaluation protocol} with verified zero drug overlap, anchor-stratified analysis, per-protein metric distributions, and cosine similarity ablation.
    \item \emph{Ablation analyses} showing that anchor quality and encoder capacity are partially substitutable, that structure-aware embeddings (ProstT5) improve transfer, and that the model learns drug-specific patterns beyond raw protein similarity.
\end{enumerate}
